\subsection{Đề 2}
\Opensolutionfile{ans}[ans/ans-0-B15]
\caulc
%%%============EX_1==============%%%
\begin{ex}%[2H5N2-2]
	Trong không gian $Oxyz$, cho đường thẳng $d\colon \dfrac{x-3}{2}=\dfrac{y-4}{-5}=\dfrac{z+1}{3}$. Véc-tơ nào dưới đây là một véc-tơ chỉ phương của $d$?
	\choice
	{$\overrightarrow{u}_2=(2; 4;-1)$}
	{\True $\overrightarrow{u}_1=(2;-5; 3)$}
	{$\overrightarrow{u}_3=(2; 5; 3)$}
	{$\overrightarrow{u}_4=(3; 4; 1)$}
	\loigiai{Véc-tơ chỉ phương của $d$ là $\overrightarrow{u}_1=(2;-5; 3)$.
	}
\end{ex}
%%%%============EX_2==============%%%
\begin{ex}%[2H5N2-2]
	Trong không gian $Oxyz$, đường thẳng $d\colon \heva{&x=2-t \\&y=1+2 t \\&z=3+t}$ có một véc-tơ chỉ phương là
	\choice
	{$\overrightarrow{u}_1=(-1; 2; 3)$}
	{$\overrightarrow{u}_3=(2; 1; 3)$}
	{\True $\overrightarrow{u}_4=(-1; 2; 1)$}
	{$\overrightarrow{u}_2=(2; 1; 1)$}
	\loigiai{$d\colon \heva{&x=2-t \\&y=1+2 t \\&z=3+t}$ có một véc-tơ chỉ phương là $\overrightarrow{u}_4=(-1; 2; 1)$.}
\end{ex}
%%%%============EX_3==============%%%
\begin{ex}%[2H5H2-2]
	Trong không gian với hệ toạ độ $Oxyz$, cho hai điểm $A(1; 1; 0)$ và $B(0;1;2)$. Véc-tơ nào dưới đây là một véc-tơ chỉ phương của đường thẳng $AB$
	\choice
	{$\overrightarrow{d}=(-1; 1; 2)$}
	{$\overrightarrow{a}=(-1; 0;-2)$}
	{\True $\overrightarrow{b}=(-1; 0; 2)$}
	{$\overrightarrow{c}=(1; 2; 2)$}
	\loigiai{
		Ta có $\overrightarrow{AB}=(-1; 0; 2)$ suy ra đường thẳng $AB$ có vectơ chỉ phương là $\overrightarrow{b}=(-1; 0; 2)$.}
\end{ex}
%%%%============EX_4==============%%%
\begin{ex}%[2H5H2-2]
	Trong không gian với hệ toạ độ $Oxyz$, cho điểm $M(1; 2; 3)$. Gọi $M_1$, $M_2$ lần lượt là hình chiếu vuông góc của $M$ lên các trục $Ox$, $Oy$. Véc-tơ nào dưới đây là một véc-tơ chỉ phương của đường thẳng $M_1M_2$
	\choice
	{$\overrightarrow{u}_4=(-1; 0; 2)$}
	{$\overrightarrow{u}_1=(0; 2; 0)$}
	{\True $\overrightarrow{u}_2=(-1; 2; 0)$}
	{$\overrightarrow{u}_3=(1; 0; 0)$}
	\loigiai{
	$M_1$ là hình chiếu của $M$ lên trục $Ox\Rightarrow M_1(1;0;0)$.\\
	$M_2$ là hình chiếu của $M$ lên trục $Oy\Rightarrow M_2(0;2;0)$.\\
	Khi đó, $\overrightarrow{M_1M_2}=(-1;2;0)$ là một véc-tơ chỉ phương của $M_1M_2$.}
\end{ex}

%%%============EX_1==============%%%
\begin{ex}%[2H5H2-3]
	Trong không gian $Oxyz$, cho hai điểm $M(1; 2; 1)$ và $N(3; 1;-2)$. Đường thẳng $MN$ có phương trình là
	\choice
	{$\dfrac{x+1}{4}=\dfrac{y+2}{3}=\dfrac{z+1}{-1}$}
	{\True $\dfrac{x-1}{2}=\dfrac{y-2}{-1}=\dfrac{z-1}{-3}$}
	{$\dfrac{x-1}{4}=\dfrac{y-2}{3}=\dfrac{z-1}{-1}$}
	{$\dfrac{x+1}{2}=\dfrac{y+2}{-1}=\dfrac{z+1}{-3}$}
	\loigiai{
	Ta có $\overrightarrow{MN}=(2;-1;-3)$.\\
	Đường thẳng $MN$ đi qua điểm $M(1; 2; 1)$ và nhận véc-tơ $\overrightarrow{MN}=(2;-1;-3)$ làm véc-tơ chỉ phương có phương trình là $\dfrac{x-1}{2}=\dfrac{y-2}{-1}=\dfrac{z-1}{-3}$.}
\end{ex}
%%%%============EX_2==============%%%
\begin{ex}%[2H5H2-3]
	Trong không gian $Oxyz$, cho hai điểm $M(1; 0; 1)$ và $N(3; 2;-1)$. Đường thẳng $MN$ có phương trình tham số là
	\choice
	{$\heva{&x=1+2t \\&y=2 t \\&z=1+t}$}
	{$\heva{&x=1+t \\&y=t \\&z=1+t}$}
	{$\heva{&x=1-t \\&y=t \\&z=1+t}$}
	{\True $\heva{&x=1+t \\&y=t \\&z=1-t}$}
	\loigiai{
		Đường thẳng $MN$ nhận $\overrightarrow{MN}=(2; 2;-2)$ hoặc $\overrightarrow{u}=(1; 1;-1)$ là véc-tơ chỉ phương.\\
		Đường thẳng $MN$ đi qua điểm điểm $M(1; 0; 1)$ và nhận $\overrightarrow{u}=(1; 1;-1)$ là véc-tơ chỉ phương có phương trình tham số là $\heva{&x=1+t \\&y=t \\&z=1-t.}$}
\end{ex}
%%%%============EX_3==============%%%
\begin{ex}%[2H5N2-3]
	Trong không gian tọa độ $Oxyz$, phương trình nào dưới đây là phương trình chính tắc của đường thẳng $d\colon\heva{&x=1+2t \\&y=3t \\&z=-2+t}(t\in \mathbb{R})$?
	\choice
	{$\dfrac{x+1}{2}=\dfrac{y}{3}=\dfrac{z-2}{1}$}
	{$\dfrac{x-1}{1}=\dfrac{y}{3}=\dfrac{z+2}{-2}$}
	{$\dfrac{x+1}{2}=\dfrac{y}{3}=\dfrac{z-2}{-2}$}
	{\True $\dfrac{x-1}{2}=\dfrac{y}{3}=\dfrac{z+2}{1}$}
	\loigiai{
		Do đường thẳng $d\colon \heva{&x=1+2t \\&y=3t \\&z=-2+t}$ đi qua điểm $M(1; 0;-2)$ và có véc-tơ chỉ phương\\ $\overrightarrow{u}=(2; 3; 1)$ nên có phương trình chính tắc là $\dfrac{x-1}{2}=\dfrac{y}{3}=\dfrac{z+2}{1}$.}
\end{ex}
%%%%============EX_4==============%%%
\begin{ex}%[2H5H2-3]
	Trong không gian với hệ tọa độ $Oxyz$, cho hai điểm $M(1;-2; 1)$, $N(0; 1; 3)$. Phương trình đường thẳng qua hai điểm $M$, $N$ là
	\choice
	{$\dfrac{x+1}{-1}=\dfrac{y-2}{3}=\dfrac{z+1}{2}$}
	{$\dfrac{x+1}{1}=\dfrac{y-3}{-2}=\dfrac{z-2}{1}$}
	{\True $\dfrac{x}{-1}=\dfrac{y-1}{3}=\dfrac{z-3}{2}$}
	{$\dfrac{x}{1}=\dfrac{y-1}{-2}=\dfrac{z-3}{1}$}
	\loigiai{
		Ta có $\overrightarrow{MN}=(-1; 3; 2)$.\\
		Đường thẳng $MN$ đi qua $N(0; 1; 3)$ và nhận $\overrightarrow{MN}=(-1; 3; 2)$ làm véc-tơ chỉ phương có phương trình là
		$\dfrac{x}{-1}=\dfrac{y-1}{3}=\dfrac{z-3}{2}$.}
\end{ex}
%%%%============EX_5==============%%%
\begin{ex}%[2H5H2-3]
	Trong không gian $Oxyz$, cho đường thẳng $d$ đi qua điểm $M(3;-1; 4)$ và có một véc-tơ chỉ phương $\overrightarrow{u}=(-2; 4; 5)$. Phương trình của $d$ là
	\choice
	{$\heva{&x=-2+3t \\&y=4-t \\&z=5+4t}$}
	{$\heva{&x=3+2t \\&y=-1+4t \\&z=4+5t}$}
	{$\heva{&x=3-2t \\&y=1+4t \\&z=4+5t}$}
	{\True $\heva{&x=3-2t \\&y=-1+4t \\&z=4+5t}$}
	\loigiai{
	Đường thẳng $d$ đi qua điểm $M(3;-1; 4)$ và có một véc-tơ chỉ phương $\overrightarrow{u}=(-2; 4; 5)$. Phương trình của $d$ là $\heva{&x=3-2 t \\ &y=-1+4t \\&z=4+5t}(t\in\mathbb{R})$.}
\end{ex}
%%%%============EX_6==============%%%
\begin{ex}%[2H5H2-3]
	Trong không gian $Oxyz$, đường thẳng $Oy$ có phương trình tham số là
	\choice
	{$\heva{&x=t\\&y=t\\&z=t}(t\in \mathbb{R})$}
	{\True $\heva{&x=0\\&y=2+t\\&z=0}(t\in \mathbb{R})$}
	{$\heva{&x=0\\&y=0\\&z=t}(t\in \mathbb{R})$}
	{$\heva{&x=t\\&y=0\\&z=0}(t\in \mathbb{R})$}
	\loigiai{Đường thẳng $Oy$ đi qua điểm $A(0;2;0)$ và nhận véc-tơ đơn vị $\overrightarrow{j}=(0;1;0)$ làm véc-tơ chỉ phương nên có phương trình tham số là $\heva{&x=0+0\cdot t\\&y=2+1\cdot t\\&z=0+0\cdot t}(t\in \mathbb{R})\Leftrightarrow \heva{&x=0\\&y=2+t\\&z=0}(t\in\mathbb{R})$.
	
}
\end{ex}

%%%============EX_1==============%%%
\begin{ex}%[2H5H2-3]
	Trong không gian $Oxyz$, cho điểm $A(1; 2; 3)$ và đường thẳng $d\colon \dfrac{x-3}{2}=\dfrac{y-1}{1}=\dfrac{z+7}{-2}$. Đường thẳng đi qua $A$, vuông góc với $d$ và cắt trục $Ox$ có phương trình là
	\choice
	{$\heva{&x=-1+2t \\&y=-2t \\&z=t}$}
	{$\heva{&x=1+t \\&y=2+2t \\&z=3+3t}$}
	{\True $\heva{&x=-1+2t \\&y=2t \\&z=3t}$}
	{$\heva{&x=1+t \\&y=2+2t \\&z=3+2t}$}
	\loigiai{
		Gọi $\Delta$ là đường thẳng cần tìm.\\
		Gọi $M=\Delta \cap Ox$. Suy ra $M\ (a; 0; 0)$.\\
		$\overrightarrow{AM}=(a-1;-2;-3)$.\\
		Đường thẳng $d$ có véc-tơ chỉ phương $\overrightarrow{u}_d=(2; 1;-2)$.\\
		Vì $\Delta \perp d$ nên $\overrightarrow{AM}\cdot \overrightarrow{u}_d=0 \Leftrightarrow 2a-2-2+6=0 \Leftrightarrow a=-1$.\\
		Vậy $\Delta$ qua $M(-1; 0; 0)$ và có véc-tơ chỉ phương $\overrightarrow{AM}=(-2;-2;-3)=-(2; 2; 3)$ nên $\Delta$ có phương trình: $\heva{&x=-1+2t \\&y=2t \\&z=3t}(t\in \mathbb{R})$.
	}
\end{ex}
%%%============EX_2==============%%%
\begin{ex}%[2H5H2-3]
	Trong không gian $Oxyz$, cho các điểm $A(1; 0; 2)$, $B(1; 2; 1)$, $C(3; 2; 0)$ và $D(1; 1; 3)$. Đường thẳng đi qua $A$ và vuông góc với mặt phẳng $(BCD)$ có phương trình là
	\choice
	{$\heva{&x=1-t \\&y=4t \\&z=2+2t}$}
	{$\heva{&x=1+t \\&y=4 \\&z=2+2t}$}
	{\True $\heva{&x=2+t \\&y=4+4t \\&z=4+2t}$}
	{$\heva{&x=1-t \\&y=2-4t \\&z=2-2t}$}
\loigiai{
Đường thẳng đi qua $A$ và vuông góc với mặt phẳng $(BCD)$ nhận véc-tơ pháp tuyến của $(BCD)$ là véc-tơ chỉ phương.\\
Ta có $\overrightarrow{BC}=(2; 0;-1)$, $\overrightarrow{BD}=(0;-1; 2)$.\\
$\Rightarrow \overrightarrow{u}_d=\overrightarrow{n}_{\left( BCD\right) }=\left[ \overrightarrow{BC}; \overrightarrow{BD}\right] =(-1;-4;-2)$.\\
Thay điểm $A(1; 0; 2)$ vào phương trình $\heva{&x=2+t \\&y=4+4t \\&z=4+2t}$ ta có $\heva{&1=2+t \\ &0=4+4 t \\&2=4+2 t} \Leftrightarrow t=-1.$\\
Vậy đường thẳng đi qua $A$ và vuông góc với mặt phẳng $(BCD)$ có phương trình là $\heva{&x=2+t \\&y=4+4t \\&z=4+2t.}$
}
\end{ex}


\Closesolutionfile{ans}
\indapan{6}{ans/ans-0-B15}

\cauds
\Opensolutionfile{ans}[ans/ans-0-B15-DS]
\begin{ex}%[2H5V2-3]
	Trong không gian $Oxyz$, cho mặt phẳng $(\alpha)\colon x+y-z+6=0$ và đường thẳng $d\colon \dfrac{x-1}{2}=\dfrac{y+4}{3}=\dfrac{z}{5}$.
	\choiceTF
	{\True Mặt phẳng $(\alpha)$ có véc-tơ pháp tuyến $\overrightarrow{n}=(1;1;-1)$}
	{\True $d\parallel (\alpha)$}
	{Đường thẳng đi qua $A$ và vuông góc với $(\alpha)$ có phương trình $\heva{&x=1+t\\&y=-4+t\\&z=t}(t\in\mathbb{R})$}
	{Hình chiếu vuông góc của $d$ lên $(\alpha)$ có phương trình là $\dfrac{x}{2}=\dfrac{y+5}{3}=\dfrac{z+1}{5}$}
	\loigiai{
		\begin{itemchoice}
			\itemch Đúng. Mặt phẳng $(\alpha)\colon x+y-z+6=0$ có véc-tơ pháp tuyến $\overrightarrow{n}=(1;1;-1)$.
			\itemch Đúng.\\ Mặt phẳng $(\alpha)\colon x+y-z+6=0$ có véc-tơ pháp tuyến $\overrightarrow{n}=(1;1;-1)$. \\Đường thẳng $d\colon \dfrac{x-1}{2}=\dfrac{y+4}{3}=\dfrac{z}{5}$ có véc-tơ chỉ phương $\overrightarrow{u}=(2;3;5)$.\\ Vì $\overrightarrow{n}\cdot\overrightarrow{u}=1\cdot2+1\cdot3-1\cdot5=0$ nên $d\parallel (\alpha)$.
			\itemch Sai. \\
			Ta có $(\alpha)\colon x+y-z+6=0$ nên $\overrightarrow{n}_{\alpha}=\left(1;1;-1 \right)$.\\
			Lấy $A(1;-4;0)\in d$.\\  Gọi $\Delta$ là đường thẳng đi qua $A(1;-4;0)$ và vuông góc với $(\alpha)$ nên $\overrightarrow{u}_{\Delta}=\left(1;1;-1 \right) $.\\ Suy ra phương trình tham số của đường thẳng $\Delta$ là $\heva{&x=1+t\\&y=-4+t\\&z=-t}(t\in \mathbb{R})$.
			\itemch Sai. \\
			Đường thẳng $d\colon \dfrac{x-1}{2}=\dfrac{y+4}{3}=\dfrac{z}{5}$ có $\overrightarrow{u}_d=(2;3;5)$.\\
			Gọi $d'$ là hình chiếu vuông góc của $d$ trên $(\alpha)\Rightarrow d'\parallel d$.\\
			Vì $d'\parallel d$ nên $d'$ nhận $\overrightarrow{u}_{d'}=\overrightarrow{u}_{\Delta}=\left(2;3;5\right)$ làm véc-tơ chỉ phương.\\
			Gọi $A'$ là hình chiếu của $A$ lên $(\alpha)$ thì $A'=\Delta \cap (\alpha)\Rightarrow A'(0;-5;1)$.\\Đường thẳng $d'$ đi qua $A'(0;-5;1)$, nhận $\overrightarrow{u}_{d'}=(2;3;5)$ làm véc-tơ chỉ phương có phương trình là $\dfrac{x}{2}=\dfrac{y+5}{3}=\dfrac{z-1}{5}$.
		\end{itemchoice}
	}
\end{ex}

\begin{ex}%[2H5H2-3]
	Trong không gian $Oxyz$, cho $M(1;0;1)$ và đường thẳng $(\Delta)\colon\dfrac{x}{1}=\dfrac{y}{2}=\dfrac{z}{3}$
	\choiceTF
	{$M\in (\Delta)$}
	{\True $(\Delta)$ có véc-tơ chỉ phương $\overrightarrow{u}=(1;2;3)$}
	{\True Phương trình tham số của $(\Delta)\colon \heva{&x=t\\&y=2t\\&z=3t}(t\in \mathbb{R})$}
	{Toạ độ hình chiếu của $M$ lên đường thẳng $(\Delta)$ là $N\left(2;4;6\right)$}
	\loigiai{
		\begin{itemchoice}
			\itemch Sai. Ta có $\dfrac{1}{1}\neq \dfrac{0}{2}\neq \dfrac{1}{3}$ nên $M\notin (\Delta)$.
			\itemch Đúng. $(\Delta)\colon\dfrac{x}{1}=\dfrac{y}{2}=\dfrac{z}{3}$  có véc-tơ chỉ phương $\overrightarrow{u}=(1;2;3)$.
			\itemch Đúng. $(\Delta)\colon\dfrac{x}{1}=\dfrac{y}{2}=\dfrac{z}{3}$ đi qua điểm $O\left(0;0;0 \right)$ và có véc-tơ chỉ phương $\overrightarrow{u}=(1;2;3)$ nên có phương trình tham số của $(\Delta)\colon \heva{&x=t\\&y=2t\\&z=3t}(t\in \mathbb{R})$.
			\itemch Sai.\\ Gọi $N(t;2t;3t)\in \Delta$ là hình chiếu vuông góc của $M$ lên $\Delta$, khi đó \allowdisplaybreaks\begin{eqnarray*}
				&&\overrightarrow{MN}\cdot \overrightarrow{u}=0\\
				&\Leftrightarrow&(t-1)+(2t-0)\cdot2+(3t-1)\cdot3=0\\
				&\Leftrightarrow&14t-4=0\\
				&\Leftrightarrow& t=\dfrac{2}{7}\\
				&\Rightarrow& N\left( \dfrac{2}{7};\dfrac{4}{7};\dfrac{6}{7}\right).
			\end{eqnarray*}
		\end{itemchoice}
	}
\end{ex}

\begin{ex}%[2H5H2-3]
	Trong không gian $Oxyz$, cho đường thẳng $d\colon \dfrac{x+1}{1}=\dfrac{y+3}{2}=\dfrac{z+2}{2}$ và điểm $A(3 ; 2 ; 0)$. 
	\choiceTF
	{\True $M(-1;-3;-2)\in d$}
	{\True $\overrightarrow{u}=(1;2;2)$ là véc-tơ chỉ phương của $d$}
	{Toạ độ hình chiếu của $A$ lên đường thẳng $d$ là $H(2;2;1)$}
	{Điểm đối xứng của điểm $A$ qua đường thẳng $d$ có tọa độ là $A'(1;0;4)$}
	\loigiai{
		\begin{itemchoice}
			\itemch Đúng. Đường thẳng $d\colon \dfrac{x+1}{1}=\dfrac{y+3}{2}=\dfrac{z+2}{2}$ nên $M(-1;-3;-2)\in d$.
			\itemch Đúng. Đường thẳng $d\colon \dfrac{x+1}{1}=\dfrac{y+3}{2}=\dfrac{z+2}{2}$ nên $\overrightarrow{u}=(1;2;2)$ là véc-tơ chỉ phương của $d$.
			\itemch Sai. \\Gọi $(P)$ là mặt phẳng đi qua $A$ và vuông góc với đường thẳng $d$.\\ Phương trình của mặt phẳng
			$(P)$ là \allowdisplaybreaks\begin{eqnarray*}
				1\cdot(x-3)+2\cdot(y-2)+2\cdot(z-0)=0 \Leftrightarrow x+2y+2z-7=0.
			\end{eqnarray*}
			Gọi $H$ là hình chiếu của $A$ lên đường thẳng $d$, khi đó $H=d \cap(P)$.\\
			Suy ra $H \in d \Rightarrow H(-1+t ;-3+2 t ;-2+2 t)$.\\ Mặt khác $H \in(P) \Rightarrow-1+t-6+4 t-4+4 t-7=0\Rightarrow t=2$. Vậy $H(1 ; 1 ; 2)$.
			\itemch Sai. \\Gọi $A'$ là điểm đối xứng với $A$ qua đường thẳng $d$, khi đó $H$ là trung điểm của $AA'$ nên ta có hệ phương trình
			\begin{eqnarray*}
				\heva{&x_A+x_A'=2x_H\\&y_A+y_A'=2y_H\\&z_A+z_A'=2z_H}\Leftrightarrow\heva{&x_A'=2x_H-x_A=2\cdot1-3=-1\\&y_A'=2y_H-y_A=2\cdot1-2=0\\&z_A'=2z_H-z_A=2\cdot2-0=4.}
			\end{eqnarray*} Do đó, $A'(-1 ; 0 ; 4)$.
		\end{itemchoice}
	}
\end{ex}

\begin{ex}%[2H5H2-3]
	Trong không gian $Oxyz$, cho mặt phẳng $(P)\colon 6x-2y+z-35=0$ và điểm $A(-1;3;6)$. 
	\choiceTF
	{\True $\overrightarrow{n}=(6;-2;-1)$ là véc-tơ pháp tuyến của $(P)$}
	{\True $A(-1;3;6)\notin (P)$}
	{\True Toạ độ điểm đối xứng với $A$ qua $(P)$ là $A'(11;-1;8)$}
	{Độ dài $OA'=186$}
	\loigiai{
		\begin{itemchoice}
			\itemch Đúng. Mặt phẳng $(P)\colon 6x-2y+z-35=0$ có $\overrightarrow{n}=(6;-2;-1)$ là véc-tơ pháp tuyến.
			\itemch Đúng. Ta có $6\cdot(-1)-2\cdot3+6-35=-41\neq 0$ nên $A\notin (P)$.
			\itemch Đúng. \\ $A'$ đối xứng với $A$ qua $(P)$ nên $AA'$ vuông góc với $(P)$.\\
			Suy ra phương trình đường thẳng $AA'\colon\heva{&x=-1+6 t \\&y=3-2 t \\&z=6+t}(t\in \mathbb{R})$.\\
			Gọi $H$ là giao điểm của $AA'$ và mặt phẳng $(P) \Rightarrow H(-1+6t;3-2t;6+t)$.
			\allowdisplaybreaks\begin{eqnarray*}
		\text{ Do }	H \in (P)&\Rightarrow& 6\cdot(-1+6 t)-2\cdot(3-2 t)+1\cdot(6+t)-35=0\\
			&\Rightarrow& 41 t-41=0\\
			&\Rightarrow& t=1 \\
			&\Rightarrow& H(5 ; 1 ; 7).
			\end{eqnarray*}
			$A'$ đối xứng với $A$ qua $(P)$ nên $H$ là trung điểm của $AA'$. Do đó, ta có hệ phương trình
			\begin{eqnarray*}
				\heva{&x_A+x_A'=2x_H\\&y_A+y_A'=2y_H\\&z_A+z_A'=2z_H}\Leftrightarrow\heva{&x_A'=2x_H-x_A=2\cdot5+1=11\\&y_A'=2y_H-y_A=2\cdot1-3=-1\\&z_A'=2z_H-z_A=2\cdot7-6=8.}
			\end{eqnarray*}
			Do đó, $A'(11 ;-1 ; 8)$.
			\itemch Sai.  Ta có $A'(11 ;-1 ; 8)\Rightarrow O A'=\sqrt{11^2+(-1)^2+8^2}=\sqrt{186}$.
		\end{itemchoice}
	}
\end{ex}
\Closesolutionfile{ans}
\indapan{3}{ans/ans-0-B15-DS}

\Opensolutionfile{ans}[ans/ans-0-B15-KQ]
\caukq
\begin{ex}%[2H5H2-6]
	Trong không gian $Oxyz$, khoảng cách từ $M(2;-4;-1)$ tới đường thẳng $\Delta \colon \heva{&x=t\\&y=2-t\\&z=3+2t}$ bằng bao nhiểu? (\textit{kết quả làm tròn đến hàng phần trăm}).
	\shortans{$7{,}48$}
	\loigiai{
	Đường thẳng $\Delta$ đi qua $N(0;2;3)$, có véc-tơ chỉ phương $\overrightarrow{u}=(1;-1;2)$.\\
	$\overrightarrow{MN}=(-2;6;4)$.\\ $\left[\overrightarrow{MN},\overrightarrow{u}\right]=\left(16;8;-4\right)$.\\
	$\mathrm{\,d}(M,\Delta)=\dfrac{\left|\left[\overrightarrow{MN},\overrightarrow{u}\right]  \right| }{\left|\overrightarrow{u} \right|  }=\dfrac{\sqrt{336}}{\sqrt{6}}=2\sqrt{14}\approx7{,}48$.
	}
\end{ex}

\begin{ex}%[2H5V2-6]
Trong không gian $Oxyz$, cho đường thẳng $d\colon \dfrac{x+1}{2}=\dfrac{y}{1}=\dfrac{z-2}{-1}$ và hai điểm $A(-1 ; 3 ; 1)$; $B(0 ; 2 ;-1)$. Gọi $C(m ; n ; p)$ là điểm thuộc đường thẳng $d$ sao cho diện tích tam giác $ABC$ bằng $2\sqrt{2}$. Tính giá trị của tổng $m+n+p$.
\shortans{$3$}
\loigiai{
Phương trình tham số của đường thẳng $d\colon\heva{&x=-1+2 t \\&y=t \\&z=2-t}(t\in \mathbb{R})$.\\
Vì $C \in d\colon\heva{&x=-1+2 t \\&y=t \\&z=2-t} \Rightarrow C(-1+2t ; t;2-t)$.\\
Ta có $\overrightarrow{AB}=(1 ;-1 ;-2)$; $\overrightarrow{AC}=(-1+2t;t;2-t) \Rightarrow\left[\overrightarrow{AB},\overrightarrow{AC}\right]=(3t-7;-3t-1; 3t-3)$.\\
Diện tích tam giác $ABC$ là $S_{ABC}=\dfrac{1}{2}\cdot\left| \left[ \overrightarrow{AB}, \overrightarrow{AC}\right] \right| =\dfrac{1}{2} \sqrt{27t^2-54t+59}$.\\
$S_{ABC}=2 \sqrt{2} \Leftrightarrow \dfrac{1}{2} \sqrt{27t^2-54t+59}=2\sqrt{2} \Leftrightarrow t=1 \Rightarrow C(1 ; 1 ; 1) \Rightarrow m+n+p=3$.}
\end{ex}

\begin{ex}%[2H5V1-5]
	Trong không gian với hệ trục tọa độ $Oxyz$. Gọi $M(a;b;c)$ thuộc đường thẳng $\Delta\colon\dfrac{x}{1}=\dfrac{y-1}{2}=\dfrac{z+2}{3}$. Biết điểm $M$ có tung độ âm và cách mặt phẳng $(Oyz)$ một khoảng bằng $2$. Xác định giá trị $T=a+b+c$.
	\shortans{$-13$}
	\loigiai{
	$M \in \Delta \Rightarrow M(t;1+2t;-2+3t)$.\\
	Ta có $\mathrm{\,d}(M ;(Oyz))=|t|=2 \Leftrightarrow\hoac{&t=2 \Rightarrow y_M=1+2t=5 \\ &t=-2 \Rightarrow y_M=1+2t=-3.}$\\
	Suy ra $t=-2$ (vì $M$ có tung độ âm). Do đó $M(-2;-3;-8)$.\\
	Vậy $a=-2$; $b=-3$; $c=-8\Rightarrow T=a+b+c=-13$.
}
\end{ex}

\begin{ex}%[2H5H2-5]
Trong không gian $Oxyz$, giao điểm của mặt phẳng $(P)\colon 3x+5y-z-2=0$ và đường thẳng $\Delta\colon \dfrac{x-12}{4}=\dfrac{y-9}{3}=\dfrac{z-1}{1}$ là điểm $M\left(x_0;y_0;z_0\right)$. Tính giá trị của tổng $x_0+y_0+z_0$.
\shortans{$-2$}
\loigiai{\allowdisplaybreaks
\begin{eqnarray*}
		&& M \in \Delta \Rightarrow M(12+4 t ; 9+3 t ; 1+t) . \\
		&& M \in(P) \Leftrightarrow 3\cdot(12+4 t)+5\cdot(9+3 t)-(1+t)-2=0 \Leftrightarrow t=-3 . \\
		&& M(0 ; 0 ;-2) \Rightarrow x_0+y_0+z_0=-2.
\end{eqnarray*}
}
\end{ex}


\begin{ex}%[2H5V2-3]
Trong không gian với hệ tọa độ $Oxyz$, cho hai điểm $A(4;6;2)$ và $B(2;-2;0)$ và mặt phẳng $(P)\colon x+y+z=0$. Xét đường thẳng $d$ thay đổi thuộc $(P)$ và đi qua $B$, gọi $H$ là hình chiếu vuông góc của $A$ trên $d$. Biết rằng khi $d$ thay đổi thì $H$ thuộc một đường tròn cố định. Tính bán kính $R$ của đường tròn đó (\textit{kết quả làm tròn đến hàng phần trăm}).
\shortans{$2{,}45$}
\loigiai{
\begin{center}
	\begin{tikzpicture}[scale=1,>=stealth, font=\footnotesize, line join=round, line cap=round]
	\def\d{7}
	\def\r{3.5}
	\path (0:0) coordinate (V)
			++(0:\d) coordinate (B')
			++(50:\r) coordinate (N)
			($(V)+(N)-(B')$) coordinate (M)	
	 ($(V)!1/6!(N)$) coordinate (K)node[above]{$d$}
	 ($(V)!4/5!(N)$) coordinate (L)
	 ($(K)!1/3!(L)$) coordinate (H)
	 ($(H)!.5!(L)$) coordinate (B)
	 ($(H)!0.5!120:(B)$) coordinate (A')
	 ($(A')!2.5!135:(H)$) coordinate (A)
		;
	\draw (V)--(B')--(N)--(M)--cycle (K)--(L) (A)--(A')--(H)--cycle (A')--(B);
	\draw pic[draw,angle radius=2mm]{right angle=B--H--A};
	\draw pic[draw,angle radius=2mm]{right angle=H--A'--A};
	\draw pic[draw,"$P$",angle radius=6mm]{angle=N--B'--V};
	\foreach \x/ \goc in {H/-30,B/-30,A'/180,A/0} 
			\fill (\x) circle (1pt)	
			($(\x)+(\goc:3mm)$) node {$\x$};
\end{tikzpicture}
\end{center}
	Dễ thấy $B \in(P)$.\\
	Gọi $A'$ là hình chiếu của $A$ trên $(P)$, $A'$ cố định.
	Ta có 
	\begin{eqnarray*}
	\heva{&AA'\perp d \\
		&AH \perp d} \Rightarrow d \perp\left(AA'H\right) \Rightarrow d \perp A'H \Rightarrow \widehat{A'HB}=90^{\circ}.
	\end{eqnarray*}
	Suy ra khi $d$ thay đổi điểm $H$ luôn thuộc đường tròn đường kính $A'B$ cố định.\\
	Trước hết ta viết phương trình đường thẳng $AA'$ qua $A$ và vuông góc với $\left( P\right) $ nên nhận
	$\overrightarrow{n}_{\left(P\right) }=\left( 1 ; 1 ; 1\right) $ là một véc-tơ phương có phương trình $\heva{&x=4+t \\ &y=6+t \\ &z=2+t.}$
	\begin{eqnarray*}
		 A'\in A A^{\prime} 
		&\Rightarrow& A^{\prime}(4+t ; 6+t ; 2+t) \\
		 A' \in(P)
		& \Rightarrow& 4+t+6+t+2+t=0 \\
		& \Leftrightarrow& 3t+12=0\\
		& \Leftrightarrow& t=-4\\
		& \Leftrightarrow& A'(0 ; 2 ;-2).
	\end{eqnarray*}
	Khi đó $A'B=\sqrt{(2-0)^2+(-2-2)^2+(0+2)^2}=2 \sqrt{6}$.\\
	Vậy $H$ luôn thuộc đường tròn cố định có bán kính $R=\dfrac{A'B}{2}=\sqrt{6}\approx 2{,}45$.
}
	
\end{ex}

\begin{ex}%[2H5V2-6]
Trong không gian với hệ tọa độ $Oxyz$, cho mặt phẳng $(P)\colon x-2y+2z-3=0$ và mặt cầu $(S)$ tâm $I(5 ;-3 ; 5)$, bán kính $R=2 \sqrt{5}$. Từ một điểm $A$ thuộc mặt phẳng $(P)$ kẻ một đường thẳng tiếp xúc với mặt cầu $(S)$ tại $B$. Tính $OA$ biết $AB=4$ (\textit{kết quả làm tròn đến hàng phần chục}).
\shortans{$3{,}3$}
\loigiai{
	Khoảng cách từ điểm $I$ đến $(P)$ là $\mathrm{\,d}(I ;(P))=\dfrac{\left| 5-2 \cdot(-3)+2 \cdot 5-3\right| }{\sqrt{1^2+(-2)^2+2^2}}=6$.\\
	$A B$ tiếp xúc với $(S)$ tại $B$ nên tam giác $A I B$ vuông tại $\mathrm{B}$, do đó ta có
	\begin{eqnarray*}
		I A=\sqrt{I B^2+A B^2}=\sqrt{R^2+A B^2}=\sqrt{\left( 2 \sqrt{5}\right)^2+4^2}=6=\mathrm{\,d}(I ;(P))
	\end{eqnarray*}
	$\Rightarrow A$ là hình chiếu của $I$ lên $(P)$.\\
	Đường thẳng $IA$ đi qua $I(5 ;-3 ; 5)$ có véc-tơ chỉ phương $\overrightarrow{u}=\overrightarrow{n}_{(P)}=(1 ;-2 ; 2)$ có phương trình $\heva{&x=5+t \\&y=-3-2t \\&z=5+2t.}(t\in \mathbb{R})$\\
	Ta có \allowdisplaybreaks\begin{eqnarray*}
		&&A=IA \cap(P)\\
		& \Rightarrow& 5+t-2\cdot(-3-2t)+2\cdot(5+2t)-3=0\\
		& \Rightarrow& t=-2 \\
		&\Rightarrow& A(3 ; 1 ; 1)\\
		& \Rightarrow& OA=\sqrt{11}\approx 3{,}3.
	\end{eqnarray*}}
	
\end{ex}
\Closesolutionfile{ans}
\indapan{6}{ans/ans-0-B15-KQ}